\begin{proof}
Assume to the contrary that such a subcategory $\mathscr{C} \supsetneq \GS_{\mathcal{E}}(T)$ exists. By \cref{conv:addsubcat}, there exists $V \in \ind(\mathscr{A})$ such that $V \in \mathscr{C}$ and $V \notin \GS_{\mathcal{E}}(T)$.
Suppose $V \in \proj(\mathscr{A})$. Since $T$ is tilting, there is a short exact sequence 
\[\begin{tikzcd}
	\xi: & 0 & V & T_0 & T_1 &  0 & 
	\arrow[from=1-2, to=1-3]
	\arrow[tail, from=1-3, to=1-4]
	\arrow[two heads,from=1-4, to=1-5]
	\arrow[from=1-5, to=1-6]
\end{tikzcd}\] with $T_0,T_1 \in \add(T)$.
Since $V,T_1 \in \mathscr{C}$ which is $\mathcal{E}$-adapted, then $\xi \in \mathcal{E}$. This yields a contradiction as $V \in \Sub_\mathcal{E}(T_0) \subseteq \GS_\mathcal{E}(T)$.

\bigskip

Assume that $V \notin \proj(\mathscr{A})$. Consider the projective resolution of $V$

\[\begin{tikzcd}
     & 0 & P_1 & P_0 & V &  0 & 
	\arrow[from=1-2, to=1-3]
	\arrow[tail, from=1-3, to=1-4]
	\arrow[two heads,from=1-4, to=1-5]
	\arrow[from=1-5, to=1-6]
\end{tikzcd}\] with $P_0,P_1 \in \proj(\mathscr{A})$. Since $T$ is tilting, there is a short exact sequence 
\[\begin{tikzcd}
     & 0 & P_1 & T_0 & T_1 &  0 & 
	\arrow[from=1-2, to=1-3]
	\arrow[tail,"\psi", from=1-3, to=1-4]
	\arrow[two heads,from=1-4, to=1-5]
	\arrow[from=1-5, to=1-6]
\end{tikzcd}\] with $T_0,T_1 \in \add(T)$.

\bigskip

We obtain the following diagram by taking the pushout along $\psi$ and by using the cokernel lifting property.
  
\[\begin{tikzcd}
         & 0 & 0 & 0 &  \\
	0 & P_1 & P_0  & V &  0 \\
     0 & T_0 & PO & V & 0 \\
      0 & T_1 & C & 0 & 0 \\
       & 0 & 0 & 0 & 
	\arrow[from=2-1, to=2-2]
	\arrow[tail, from=2-2, to=2-3]
	\arrow[two heads,from=2-3, to=2-4]
	\arrow[from=2-4, to=2-5]
        \arrow[from=3-1, to=3-2]
	\arrow[tail, from=3-2, to=3-3]
	\arrow[two heads,from=3-3, to=3-4]
	\arrow[from=3-4, to=3-5]
        \arrow[from=4-1, to=4-2]
	\arrow[tail, from=4-2, to=4-3]
	\arrow[two heads,from=4-3, to=4-4]
	\arrow[from=4-4, to=4-5]
        \arrow[from=4-2, to=5-2]
	\arrow[two heads, from=3-2, to=4-2]
	\arrow[tail,"\psi",from=2-2, to=3-2]
	\arrow[from=1-2, to=2-2]
        \arrow[from=4-3, to=5-3]
	\arrow[two heads, from=3-3, to=4-3]
	\arrow[tail,from=2-3, to=3-3]
	\arrow[from=1-3, to=2-3]
        \arrow[from=4-4, to=5-4]
	\arrow[two heads, from=3-4, to=4-4]
	\arrow[equals,from=2-4, to=3-4]
	\arrow[from=1-4, to=2-4]
\end{tikzcd} \]

As the short sequences given by the three columns and those given by the first two rows of the above diagram are exact, by \cref{lem:3x3}, the short sequence given by the last row is exact too. So $C \cong T_1$. Since $V,T_0 \in \mathscr{C}$ which is $\mathcal{E}$-adapted extension-closed, then the middle row lies in $\mathcal{E}$ and $W = PO \in \mathscr{C}$.

Since $T$ is tilting, there is also a short exact sequence 
\[\begin{tikzcd}
     & 0 & P_0 & T_2 & T_3 &  0 & 
	\arrow[from=1-2, to=1-3]
	\arrow[tail,"\phi", from=1-3, to=1-4]
	\arrow[two heads,from=1-4, to=1-5]
	\arrow[from=1-5, to=1-6]
\end{tikzcd}\] with $T_2,T_3 \in \add(T)$.


Given this short exact sequence and the middle column of the diagram above, we obtain the following diagram by taking the pushout along $\phi$ and by using the cokernel lifting property again.

\[\begin{tikzcd}
         & 0 & 0 & 0 &  \\
	0 & P_0 & T_2  & T_3 &  0 \\
     0 & W & PO' & C' & 0 \\
      0 & T_1 & T_1 & 0 & 0 \\
       & 0 & 0 & 0 & 
	\arrow[from=2-1, to=2-2]
	\arrow[tail,"\phi", from=2-2, to=2-3]
	\arrow[two heads,from=2-3, to=2-4]
	\arrow[from=2-4, to=2-5]
        \arrow[from=3-1, to=3-2]
	\arrow[tail, from=3-2, to=3-3]
	\arrow[two heads,from=3-3, to=3-4]
	\arrow[from=3-4, to=3-5]
        \arrow[from=4-1, to=4-2]
	\arrow[equals, from=4-2, to=4-3]
	\arrow[two heads,from=4-3, to=4-4]
	\arrow[from=4-4, to=4-5]
        \arrow[from=4-2, to=5-2]
	\arrow[two heads, from=3-2, to=4-2]
	\arrow[tail,from=2-2, to=3-2]
	\arrow[from=1-2, to=2-2]
        \arrow[from=4-3, to=5-3]
	\arrow[two heads, from=3-3, to=4-3]
	\arrow[tail,from=2-3, to=3-3]
	\arrow[from=1-3, to=2-3]
        \arrow[from=4-4, to=5-4]
	\arrow[two heads, from=3-4, to=4-4]
	\arrow[from=2-4, to=3-4]
	\arrow[from=1-4, to=2-4]
\end{tikzcd} \]


As the short sequences given by the three rows and those given by the first two columns of the above diagram are exact, by \cref{lem:3x3}, the short sequence given by the last column is exact too. So $C' \cong T_3$. Since $T$ is tilting, the middle column splits, hence $PO' \cong T_1 \oplus T_2$.


We get the short exact sequence 
\[\begin{tikzcd}
    \xi: & 0 & W & T_1 \oplus T_2 & T_3 &  0 & 
	\arrow[from=1-2, to=1-3]
	\arrow[tail, from=1-3, to=1-4]
	\arrow[two heads,from=1-4, to=1-5]
	\arrow[from=1-5, to=1-6]
\end{tikzcd}\]
Given that $W,T_3 \in \mathscr{C}$ which is $\mathcal{E}$-adapted, then $\xi \in \mathcal{E}$ and $W \in \Sub_\mathcal{E}(T_1 \oplus T_2) \subseteq \GS_\mathcal{E}(T)$.

We finally obtain the following short exact sequence from the middle row of the first diagram
\[\begin{tikzcd}
    \xi': & 0 & T_0 & W & V &  0 & 
	\arrow[from=1-2, to=1-3]
	\arrow[tail, from=1-3, to=1-4]
	\arrow[two heads,from=1-4, to=1-5]
	\arrow[from=1-5, to=1-6]
\end{tikzcd}\]

Since $\xi' \in \mathcal{E}$ and $V' \in \GS_\mathcal{E}(T)$, then we have a contradiction as $V \in \Gen_\mathcal{E}(W) \subseteq \GS_\mathcal{E}(T)$. 
\end{proof}

\ibenj{Je crois bien que ta preuve fonctionne superbement ! Bravo ! Et petit plus, ta preuve semble fonctionner dans un cadre en dehors du type $A$ --- vraisemblablement, pour n'importer quelle catégorie abélienne héréditaire. A checker, mais c'est déjà une très belle preuve !}
\ibenj{Eh ! Mais ça peut expliquer aussi pourquoi le calcul de $\GS_\mathcal{E}$ se fait avec au plus deux étapes non ? Je pense que ça peut être un corollaire.}

